\section{Discussion}
\label{sec:discussion}

\subsection{Split strategy against seed}

A shared scaffold means the test molecule's ring system was already measured in training, usually as an analogue differing by a methyl or a halogen --- a near-duplicate in feature space that deduplication does not remove, so a model can match on scaffold identity and score well without having learned anything transferable. The gap between the two splits is therefore not a defect to be removed; it estimates how much of the random-split score was recognition of a familiar scaffold.

Crossing three seeds with two split strategies separates the two sources of variation, and the split strategy dominates: it moves the mean by 0.0354 against seed standard deviations of 0.0022 and 0.0025, with no overlap between the two groups of runs, so the drop from 0.915 to 0.880 is caused by the protocol and not by initialisation luck. The \emph{trajectory} is nonetheless seed-sensitive --- best epoch 28 to 74 at nearly equal final quality --- which argues for checkpointing on a validation metric rather than a fixed epoch count.

\subsection{Overfitting and calibration}

The divergence in \cref{fig:curves} is unambiguous overfitting, as expected with \num{67588} parameters and 1375 training molecules; nothing underfits, since training loss is low and still falling when training stops. The subtlety is that validation ROC-AUC does not degrade as validation loss rises: the network becomes over-confident, pushing probabilities towards 0 and 1, which cross-entropy punishes severely on the molecules it gets wrong, while the \emph{ordering} of molecules by predicted probability stays stable or improves. Hence the ROC-AUC checkpoint rule, and hence the operating point must be selected on validation data rather than left at 0.5. The loss spikes and the sawtooth come from 295-molecule validation subsets and from batch-norm statistics shifting between modes; they do not grow, so this is noise around a converged model.

\subsection{Baseline against network}

The network wins the threshold-free metrics and loses the thresholded ones: it ranks molecules better while classifying them slightly worse at its own operating point, which is the same calibration problem seen from the other side. Its thresholds, $0.424 \pm 0.213$ and $0.748 \pm 0.248$, are spread so widely that the three seeds effectively use different operating points, whereas the strongly regularised linear model produces well-behaved probabilities and converts a lower ROC-AUC into higher accuracy and $F_1$.

Does the evidence support the claim that the network is better? Only weakly, and only for ranking. On the scaffold split the difference is 0.009 ROC-AUC against a seed standard deviation of 0.025, so the network's seeds overlap the baseline entirely and \textbf{the two models are indistinguishable there}; the random split's 0.020 is about 2.5 times the seed spread, suggestive but resting on three runs and one split seed. A fingerprint matching a graph network on BBBP is well replicated rather than anomalous \parencite{Chemprop}: with 1375 training molecules there is little for \num{67588} parameters to learn beyond substructure presence.

\subsection{Do the errors follow from similarity?}

The four most confident errors come from the sparsest regions of the training distribution, and similarity explains three of them. Both false positives inherit a positive neighbour: the bis-glutarimide (true 0, $p = 0.994$) matches a simple glutarimide at $T = 0.294$ labelled permeable, so the network transfers that motif to a molecule holding \emph{two} of them joined by a polar linker --- exactly the change that doubles polar surface area and blocks penetration --- and the macrocycle (true 0, $p = 0.989$) is the same story at $T = 0.220$. A shared local substructure drives the prediction while the global property that decides permeability, the size and polarity of the whole molecule, is what differs. The second false negative mirrors this: chromone-2-carboxylic acid (true 1, $p = 0.005$) matches a bis-chromone labelled non-permeable, and the network copies that label onto roughly half of it.

The first false negative is where similarity explains nothing: the cyclic acetal (true 1, $p = 0.005$) has methylcyclopentane as its nearest training molecule at $T = 0.200$, labelled permeable and predicted 1.000, and such a neighbour should pull the prediction up. The cause is distributional --- at 7 heavy atoms it lies far below the training median of 23, and mean pooling over 7 atoms of which 3 are ring oxygens yields an embedding dominated by polar-oxygen signal, \cref{eq:pool} normalising by atom count so that a small molecule has no averaging to dilute an unrepresentative atom. Low similarity does not make the model hedge either: the macrocycle sits 0.220 from anything in training and is still assigned $p = 0.989$, so predicted probability says nothing about whether a molecule lies inside the model's domain of applicability. A nearest-neighbour distance check would have flagged all four cases.

\subsection{Limitations and what the score means}

\textbf{Label uncertainty.} Ten molecules carry both labels and every record from row 1000 onwards is positive, indicating a concatenation of sources with different labelling conventions rather than one experiment; the underlying quantity is continuous and assay-dependent, and binarising it discards the margin, so the ceiling on achievable accuracy is set by this noise, not by the model. \textbf{Dataset size.} The seed standard deviation of 0.025 ROC-AUC is of the same order as the entire baseline-versus-network difference, and 296 test molecules support a confidence interval of roughly $\pm 0.04$, so no ranking between the models is established. \textbf{No external validation.} Every number comes from an internal re-split of one public file. \textbf{Narrow applicability and poor calibration.} The model is confident where it has almost no support, nothing abstains, and thresholds ranged from 0.42 to 0.95, so the outputs cannot be read as probabilities.

The majority-class predictor reaches $F_1$ 0.865 by calling everything permeable, against which the baseline's 0.916 is a real but modest gain: any report leading with accuracy or $F_1$ here is largely reporting the base rate. More fundamentally, the model encodes no mechanism --- neither passive diffusion, active transport, P-glycoprotein efflux nor plasma protein binding --- and cannot distinguish a molecule that fails to cross the barrier from one that crosses and is pumped straight back out. A high predicted probability states only that a molecule resembles compounds this file labelled permeable, not that it will reach the brain, and certainly not that any structural feature \emph{causes} permeability. Nothing here supports a clinical decision, and the error analysis shows how one would fail.
